\section{Future work}\label{S:future-work}

The second-moment argument suggests several questions about higher moments.
We discuss them here as possible extensions. The results of this manuscript
concern the second moment; the proposed limits below require further proof.
We retain Hypothesis~\ref{H:homog}, the initial data $(1,0)$ and the exponent
$\kappa$ from~\eqref{E:kappa}.

\paragraph{Higher even moments.}
A natural first question is whether, for each fixed even integer $p \ge 2$,
\begin{equation}\label{E:future-even-moments}
  \lim_{t\to\infty}t^{-\kappa}\log\E|u(t,x)|^p
  =\left(\frac{p}{2}\right)^{\frac{4-\alpha}{3-\alpha}}
  C_2(\Gamma,\alpha_0),\qquad x\in\R^d.
\end{equation}
The candidate comes from the variational expression in
Theorem~\ref{T:intro-second-moment}. If a Gaussian shift produces an integrated
deterministic growth rate $I_\Psi$, its contribution to the $p$th moment should
be compared with the same Gaussian cost $v_\Psi/2$. This suggests optimizing
$pI_\Psi-v_\Psi/2$ over the full positive-control class. The homogeneity of
these variational values is already proved for every positive coefficient in
Proposition~\ref{P:variational-positive-scores}. What remains is their
identification with the higher point-moment limits. Formula
\eqref{E:future-even-moments} uses the intrinsic second-moment constant; it does
not require identifying that constant with the one predicted in the source.

Even orders allow us to use $|u|^p = u^p$ and Gaussian pairing expansions
without assuming that the solution is nonnegative. A $p$-chain expansion must
retain self-pairings, cross-pairings, unequal chain lengths and empty chains.
The scaling identity and the Gaussian change of measure offer a different
starting point that may avoid estimating all these pairings separately.
The lower-bound construction and finite Gram optimization also suggest
extensions with coefficient $p$. Higher-moment convergence, uniform
integrability and the comparison with the original point observable must
still be justified.

\paragraph{Point moments and spatial regularity.}
The main obstacle in the present upper-bound method occurs in
Proposition~\ref{P:original-point-transfer}. Stationarity and finite
propagation identify the pointwise second moment with an expected spatial
$L^2$ norm. The displacement component of the wave phase controls this norm.
For a $p$th moment, the corresponding argument would require a spatial $L^p$
bound. Taking a higher power of an $L^2$ bound does not supply it.

The existing displacement space is $H^{1/2}(\R^d)$. On a fixed ball, the
fractional Sobolev embeddings control every finite $L^p$ norm when $d = 1$,
$L^p$ up to $p = 4$ when $d = 2$, and $L^p$ up to $p = 3$ when $d = 3$;
see~\cite{di-nezza.palatucci.ea:12:hitchhikers}*{Section~6}.
Thus, the present phase space suggests a route to all fixed even moments in
one dimension and to the fourth moment in two dimensions. In three dimensions,
the fourth moment already lies beyond this direct embedding argument.
This is a restriction of the method, not a nonexistence statement for the
moment. A different point-moment comparison or stronger spatial estimates
would be needed. Any additional regularity must be established under the
original covariance assumptions.

\paragraph{A first test and further questions.}
The fourth moment in one spatial dimension is a useful first test. One would
first establish the fourth-moment point comparison and the required source
limits. The Gaussian upper argument would then need to control the fourth
power of the evolution norm, with bounds uniform in the growing spatial
radius. Matching this bound with the shifted lower estimate would test the
prediction $C_4 = 2^{(4-\alpha)/(3-\alpha)}C_2$. The validity of the Sobolev
embedding alone does not complete these steps.

Even a proof for every fixed even $p$ would leave two separate questions:
absolute moments of other orders, and estimates uniform as $p\to\infty$.
Interpolation between adjacent even moments need not identify the exact
constant at an intermediate order. Uniform large-$p$ estimates also require
tracking the dependence on $p$ throughout the construction, source limits and
Gaussian tail bounds. These questions may clarify which parts of the present
method apply to moment growth beyond the second moment.
